A continuación se detallan las funcionalidades disponibles en los dos módulos principales del sistema.

\paragraph{Feed the Realm - Game}
\begin{itemize}
  \item Exploración y unión a mundos creados por la comunidad mediante un feed global de mundos.
  \item Creación y personalización de personaje con nombre, biografía y apariencia únicos.
  \item Persistencia de progreso, inventario y logros del personaje entre sesiones.
  \item Desplazamiento entre zonas y áreas conectadas mediante portales.
  \item Multijugador en tiempo real con personajes visibles, nametags y chat flotante.
  \item Aceptación, seguimiento y completado de quests otorgadas por NPCs, con múltiples tipos de objetivos.
  \item Combate contra enemigos con armas de cuerpo a cuerpo, a distancia y de área de efecto.
  \item Gestión de \textit{stamina}, inventario, equipamiento, consumibles y botín.
  \item Recolección de ítems, oro y recompensas de enemigos derrotados y cofres.
  \item Equipamiento de armas mediante slots de acceso rápido y comparación de estadísticas de ítems.
  \item Compra y uso de consumibles, equipamiento y cosméticos.
  \item Interacción con NPCs mediante sistemas de diálogo personalizables.
  \item Intercambio de oro por ítems en tiendas y comerciantes del mundo.
  \item Compra de gemas y cosméticos premium a través del marketplace integrado.
  \item Inventario global de cosméticos compartido entre todos los mundos.
  \item Efectos visuales, música ambiental y efectos de sonido.
  \item Visualización de información detallada del mundo antes de unirse.
  \item Comportamientos inteligentes de enemigos y NPCs, incluyendo combate, persecución y patrullaje.
  \item Gestión segura de cuentas con autenticación y recuperación de acceso.
\end{itemize}

\paragraph{Feed the Realm - World Editor}
\begin{itemize}
  \item Creación y publicación de mundos personalizados para la comunidad de Feed the Realm.
  \item Organización de mundos en múltiples zonas interconectadas.
  \item Publicación de zonas individuales o mundos completos.
  \item Colocación de edificios, objetos de entorno, bases, rampas y assets decorativos.
  \item Personalización de materiales de suelo, colores y presentación visual.
  \item Definición de puntos de aparición de jugadores y áreas de spawn de enemigos.
  \item Configuración de generadores de enemigos, NPCs y cofres de botín.
  \item Importación de modelos 3D y assets personalizados para uso en los mundos.
  \item Acceso a una biblioteca de modelos predeterminados y contenido reutilizable.
  \item Creación y personalización de enemigos, NPCs y sus apariencias visuales.
  \item Configuración de árboles de diálogo de NPCs e interacciones con quests.
  \item Construcción de cadenas de quests completas con objetivos, recompensas e integración de diálogos.
  \item Creación de objetivos basados en derrotas de enemigos, recolección de ítems, conversaciones, adquisición de oro y más.
  \item Configuración de recompensas de quests incluyendo ítems y moneda.
  \item Creación de armas, consumibles, tablas de botín y drops de ítems.
  \item Asignación de tablas de drops y recompensas de botín a enemigos.
  \item Construcción de tiendas en el mundo y definición de ítems comprables.
  \item Configuración de precios de ítems en oro, gemas o ambos.
  \item Creación de portales y vinculación a zonas específicas.
  \item Carga de datos y assets del mundo directamente al servicio de publicación.
  \item Monitoreo del estado del mundo y los servidores desde el editor.
  \item Gestión de contenido del mundo mediante una interfaz mejorada con indicadores de flujo de trabajo.
  \item Compartición instantánea de mundos publicados a través del feed de descubrimiento de la comunidad.
\end{itemize}

\vspace{0.25cm}
El conjunto de funcionalidades descritas refleja un sistema completo que abarca tanto la experiencia de juego 
como las herramientas de creación de contenido. Desde la gestión de personajes y la progresión mediante quests 
hasta la publicación de mundos personalizados, Feed the Realm busca ofrecer una plataforma donde los jugadores 
y creadores puedan interactuar en un ecosistema compartido, donde el contenido generado por la comunidad es el 
eje central de la experiencia.